\section{Limitations and Mitigations}\label{sec:tradeoffs}\label{sec:mitigations}

\subsection{Excluded transactions are also disclosed}\label{sec:noninclusion}

A time lock opens whether or not its transaction is included. If a builder receives a sealed transaction and leaves it out, the builder, or anyone else holding it, can read it after the delay. A threshold committee that decrypts only included transactions keeps excluded ones private, which matters because an excluded transaction can still reveal an intent to trade. The secrecy \leanseal gives is temporary in every case: it suits protection against front-running, not the lasting confidentiality that shielded transfers need.

This is not specific to \leanseal. Call a fallback opening \emph{input-free} if it takes only the sealed transaction and auxiliary data that is available whether or not the transaction is included, with no later input or secret release that depends on inclusion.

\begin{claim}[informal]\label{claim:noninclusion}
In any scheme whose fallback opening is input-free, whoever can open an included sealed transaction without the sender can open it equally well when it is not included.
\end{claim}

Since inclusion is not an input to the opening, the opening cannot depend on it. Revealing a transaction if and only if it is included therefore needs either a party that can refuse to open, such as a committee, a TEE or the sender, or an opening that takes a later inclusion witness as input, as witness encryption with a blockchain clock does~\cite{LJKW18}. The second class is outside the claim; we know of no practical construction in it, and the later input must come from some party or from consensus. Naraghi draws the same three-way distinction among the inputs of a forced opening~\cite{Naraghi26}. This disclosure is inherent in the guarantee: a puzzle protects against a withholding sender only because the sender has released it before the ordering is fixed. Four measures delay the disclosure, narrow its audience and reduce its value.
\begin{itemize}[leftmargin=1.5em]
  \item \textbf{Delay.} An excluded transaction opens a delay after exposure, not at exposure as in today's public mempool; by then, most opportunities to front-run it have passed.
  \item \textbf{Expiry.} A short validity window in the signed transaction means that by the time an excluded transaction can be read, it can no longer be executed. What leaks is the intent. The sender can also cancel early by spending the nonce.
  \item \textbf{Fewer holders.} Sent only to the parties that can include it in the current slot, the builder or the inclusion-list committee, rather than gossiped, an excluded transaction is readable only by them, and exposure starts close to $\torder$. Broadcasting only a hash of the sealed transaction and publishing the rest once its slot is reserved~\cite{Herold26,Naraghi26} also starts exposure close to $\torder$, and a transaction whose commitment is never included is never exposed; but the sender can then decline to publish, which the deposit must also cover (\cref{sec:withholding}).
  \item \textbf{A committee.} Encrypting under a key derived from both a \leanseal key and a committee's key keeps excluded transactions private while the committee is honest (\cref{sec:composition}).
\end{itemize}

\subsection{Withholding: disclosure is forced, execution is not}\label{sec:withholding}

With plain commit--reveal, the sender holds a free option: it reveals its key if the ordering is favorable and withholds it otherwise. The time lock turns withholding into delay, so withholding cannot keep a well-formed puzzle's transaction hidden. Execution, however, cannot wait at the committed position for minutes of solving without stalling every later transaction. A late-opened transaction must therefore execute later or be dropped, and with a tight deadline, a late execution that reverts amounts to a cancellation. Executing it later is safe only if no one could read its key earlier: a builder that keeps a pending reveal out until the deadline holds the key, and could trade around the late execution, so our prototype drops it (\cref{sec:poc}). Withholding thus still lets the sender delay or cancel execution, and a deposit must deter that choice. The reveal deadline makes a missed reveal observable; the time lock separately guarantees disclosure.
\begin{itemize}[leftmargin=1.5em]
  \item \textbf{A deposit forfeited on a missed reveal.} A missed reveal is visible on chain. The deposit is a plaintext field, since the protocol cannot see the transaction's value; LUCID's partly refunded top-of-block fee~\cite{LUCID} is one such rule. The penalty must be irrecoverable: the sender knows its own key, so any part of the deposit paid to whoever publishes the key is a part the sender can claim back. A forfeitable fee is also much weaker against an attacker with two or more commitments ahead of its victim (fee inversion~\cite{Naraghi26}, \cref{sec:comparison}). Capping at one the commitments that may execute ahead of a transaction, or making release decisions blind (below), avoids this.
  \item \textbf{The value of the option.} During the reveal window the sender sees the ordering and the market, but not the other sealed transactions. Its option is roughly an at-the-money option over the window, worth about $0.4\,\sigma\sqrt{\Delta t}$ of the traded value: about 0.015\%, or \$1.50 on a \$10,000 trade, for an asset with 60\% annual volatility and a 12-second window. If a late-opened transaction may still execute, the option lasts until the forced opening instead: about 0.02\% if fast solvers open the reference puzzle in 30 seconds and 0.06\% for the 3.3 minutes one laptop core takes.
  \item \textbf{Sizing the deposit.} These are values of an unpenalized option in a volatility-only model, not sufficient deposits. A deposit charged only when the option is exercised raises the exercise threshold but leaves a residual option: a deposit equal to the option's value still leaves more than half of it. The ordering itself also carries information that the model omits. Sizing the deposit therefore needs a target exercise probability or a bound on the gain from withholding~\cite{Naraghi26}, over the whole window in which the option lasts.
  \item \textbf{The last revealer.} A sender that sees the others' reveals before its own deadline decides with more information. Releasing reveals simultaneously is an additional protocol requirement; a party that merely collects plaintext keys could leak them early. Blind simultaneous release~\cite{Naraghi26} is one mechanism: right after ordering, each releaser, here each sender, escrows its key under a timed pad that outlasts a window of one or two seconds, so every decision to reveal is fixed before any key is public. It needs a committee to fix the set of timely escrows, funded fast solvers and a slashable deposit, and escrows by end users remain an open problem (\cref{sec:comparison}).
\end{itemize}
Applications in which order within a batch does not matter avoid this problem, since settlement can wait for every forced opening; \cref{sec:auction} specifies one.

\paragraph{What the puzzle adds to commit--reveal.}
With the same deposit and deadline, plain commit--reveal already makes a missed reveal observable and penalized, so in a general mempool the deposit, not the puzzle, deters withholding. The puzzle adds two things, each conditional on a well-formed puzzle. First, a withheld transaction becomes public, so withholding no longer keeps the sender's intent private: a sender cannot submit several sealed orders and reveal only the profitable ones without disclosing the others. Second, a sender that goes offline after submitting, or whose key publisher fails, does not lose its transaction: anyone can recover it, and a protocol that executes late openings, such as the auction of \cref{sec:auction}, can still run it.

\paragraph{Strategic senders.}
A sender could build a puzzle whose key matches $C$ but whose chains lead elsewhere, or encrypt its transaction under a key other than the committed one; either way, the solved key does not open the transaction. No efficient check of puzzles at submission is known (\cref{sec:open}). A malformed puzzle costs nothing to make, while a well-formed one costs about 21\,s of a laptop and removes the option to withhold privately. Without a penalty for malformation, a sender that values the option therefore never builds a well-formed puzzle, and the two additions above cover only senders whose software builds puzzles honestly: senders that go offline or whose key publisher fails, not senders that withhold by strategy. Against strategic senders the deposit alone then deters withholding, exactly as in commit--reveal. \Cref{sec:bond} adds the missing penalty.

\subsection{A bond for malformed puzzles}\label{sec:bond}

A missed reveal is observable at once; a malformed puzzle only after a forced opening fails. A protocol can charge the two separately: the deposit $d$ on a missed reveal, and a bond $B$ when malformation is established. The bond is returned when the sender's reveal is accepted or when any key matching $C$ opens the transaction, and is forfeited if malformation is established first.

\paragraph{Proving a failed opening.}
Given a key $K$ with $\Hcommit(\ctx, K) = C$, a contract checks the envelope's authentication tag directly, which needs no decryption and costs work linear in the envelope's length, which admission caps (\cref{sec:auction}). A failing tag establishes malformation at once. Chains that do not lead to the committed key are harder: the evidence is the whole computation of $\Open$, $N + O(n)$ steps, which no contract can repeat. An optimistic verification game, as in TrueBit~\cite{TR17} and Arbitrum~\cite{Arbitrum18}, can establish it in about $\log_2 N$ rounds. We outline one below; a deployment needs a full specification and its gas costs (\cref{sec:open}).

\paragraph{The game.}
A \emph{state} of $\Open$ is the solver's checkpoint (\texttt{SPEC.md}, Section~6.4: the chain index, the steps done in that chain, the chain state and the accumulator) together with the current link, a few hundred bytes. The contract computes the hash of the initial state from the puzzle itself. A claimant stakes and posts the final state, whose key $K'$ the contract checks to satisfy $\Hcommit(\ctx, K') \neq C$. Within a challenge period, anyone may stake and dispute the claim. Each challenger then plays its own game against the claimant: the claimant posts the hash of the state at the midpoint of the interval still in dispute, and the challenger says whether it agrees, which halves the interval.

When one step remains, the claimant opens the state both sides agreed on, which for the first step is the initial state. The contract recomputes its successor, one SHA-256 compression or, at a chain boundary, a few labelled hashes, and compares the result's hash with the claimant's hash for the end of the interval: the claimant wins that game if they match, and the challenger otherwise. A party that misses a move's deadline loses the game, and the loser forfeits its stake. For the reference puzzle, $N + O(n)$ is just above $2^{33}$, so a game takes 34 rounds of two moves, each a transaction carrying at most one 32-byte hash, and one final step. Naraghi assumes a similar mechanism, without analysis, for the pads of blind simultaneous release~\cite{Naraghi26}.

\paragraph{Who decides.}
A claim stands only if no one challenges it within the period, or if the claimant wins every game. Games run in parallel, so a challenger colluding with the claimant that loses on purpose cannot defeat an honest one. Since $\Open$ is deterministic and its input public, anyone who has solved a well-formed puzzle can challenge a false claim, so a sender that went offline needs no one to act for it. This is the usual assumption of optimistic verification: at least one honest party solves the puzzle within the challenge period and gets each move included before its deadline. A key that matches $C$ does not refute a claim, since a malicious sender knows such a key; posted before the claim stands, it opens the transaction, ends the games, returns every stake and returns the bond.

\paragraph{Paying solvers.}
A sender that withholds its key can still publish it after the deadline, so a reward for the first matching key would go to the sender and leave the solvers that did the work unpaid, and a solver of a malformed puzzle has no key to post.

The protocol therefore assigns each missed reveal to a solver drawn from a staked registry. Within a deadline $D_s$, the solver must submit a key matching $C$, evidence that the tag fails under such a key, or a malformation claim; a key posted by anyone also counts. A timely claim keeps the assignment until the claim is decided. The solver is paid from the forfeited deposit, and recovers its stake, for a payable outcome: a valid opening, a failing tag, or its claim standing. A solver that submits nothing in time, or whose claim fails, loses its stake, and the transaction passes to the next solver. Designation prevents outsiders from copying an assignment's reward, but does not exclude solvers that the sender controls. The effective penalty is therefore $d - r$, or $d + B - r$ when malformation is established, where $r$ is every payment the sender or its coalition can recover; a forfeited bond is burned, so that $r$ is at most the solver's payment.

\paragraph{What the bond changes.}
\Cref{tab:bond} lists what each choice forfeits. A sender that withholds by strategy now either discloses its transaction or, to keep it private, forfeits $d + B$, while one that goes offline with a well-formed puzzle forfeits only $d$. Since $B$ is charged only when malformation is established, it can be much larger than $d$ without raising the cost of an honest failure. Commit--reveal cannot make this distinction, since an accidental missed reveal and a strategic one look the same; it can raise the price of the option only by raising the price of going offline. In a general mempool where late openings do not run, as in our prototype, the bond makes concealed withholding cost $d + B$, but withholding still cancels the transaction for $d$; in an auction, a valid bid from a well-formed puzzle settles anyway (\cref{sec:auction}).

\begin{table}[t]
  \centering\small
  \caption{What a sender forfeits for each choice, with a deposit $d$ forfeited on a missed reveal and a bond $B$ forfeited when malformation is established, before subtracting payments the sender can recover. Building and locked capital are costs too, and are not shown.}
  \label{tab:bond}
  \begin{tabularx}{\textwidth}{@{}>{\raggedright\arraybackslash}p{4.1cm}>{\raggedright\arraybackslash}X>{\raggedright\arraybackslash}X@{}}
    \toprule
    \colhead{Sender} & \colhead{Commit--reveal, deposit $d$} & \colhead{\leanseal, deposit $d$ and bond $B$} \\
    \midrule
    Reveals on time & nothing & nothing \\
    Goes offline & $d$; the transaction is lost & $d$; the transaction is disclosed, and runs where late openings run \\
    Withholds, well-formed puzzle & $d$; the transaction stays private & $d$; the transaction is disclosed \\
    Withholds, malformed puzzle & n/a & $d + B$ if it stays private; $d$ if the sender later posts a key that opens it \\
    \bottomrule
  \end{tabularx}
\end{table}

The bond gives no assurance when a transaction is accepted: it prices malformation but does not rule it out, so a protocol still cannot know at submission that a puzzle will open, and a sender that values its option above $d + B$ still buys it.

\paragraph{Denial of service.}
A missed reveal can induce nearly a full solve while costing only the net deposit penalty $d - r$: a sender whose chains are malformed, but who holds a key matching $C$ and a valid envelope under it, can post the key just before the solvers finish, which opens the transaction and returns the bond. The bond is additionally forfeited only when malformation is established before a successful opening. The net deposit must therefore pay for a solve at the admission limits, together with any claim, game and redundant solver. For the reference puzzle, a solve takes about 3.3 minutes of a laptop core, or 30\,s of a circuit as fast as the assumed attacker; on rented laptop-class cores at four cents per core-hour it costs about a fifth of a cent.

Deposits price demand; they do not by themselves provide solver capacity, and capacity at attacker speed is scarcer. An attacker that exhausts it delays forced openings; it does not make any transaction readable earlier, since secrecy rests on the delay, not on the solvers. A cap on open puzzles per sender does not bound the total load when the attacker uses many accounts; admission limits and deposits do.

\subsection{A complete application: sealed-bid batch auctions}\label{sec:auction}

In a sealed-bid batch auction, order within a batch does not matter, so settlement can wait until every bid is open, and a forced opening still settles a bid at its committed terms. This is the classical use of timed commitments~\cite{BN00}, and the setting where the puzzle adds most. We specify a frequent batch auction that sells, in each slot, a fixed quantity of one asset for a quote asset at a uniform price, and compare it with the same auction under commit--reveal. Its parameters are the deposit $d$, the bond $B$, a reveal window of $W$ slots, a solver deadline $D_s$, a challenge period $D_c$ and a move deadline $\tau$.
\begin{enumerate}[leftmargin=1.5em]
  \item \textbf{Supply and bids.} Before batch $k$ opens, the seller deposits its quantity $Q_k$ in the contract, so each batch settles from its own supply and escrows. A bid is a plaintext $(k, p, x)$: the batch, a unit price and a quantity, each a 32-byte unsigned integer, 96 bytes in all. It is \emph{valid} if its batch is $k$, $p$ is at least the reserve price, $x \le Q_k$, and $p x \le E$, where $E$ is the escrow posted in plaintext with the bid.
  \item \textbf{Admission.} The contract accepts a sealed bid for batch $k$ only in the block of slot $k$, with its escrow, $d$ and $B$, and only with the reference work and fixed limits: $nm = 2^{33}$, $n \le 16$, $b = 24$, an envelope of the bid's fixed length, and at most a fixed number of bids per batch. Checking the commitment, deriving the subkey and checking the tag then take at most 17 SHA-256 compressions.
  \item \textbf{Exposure.} A bid is exposed when it first leaves the bidder, to a builder or a mempool, and the set of bids is fixed when slot $k$'s block is final. The bidder sends its bid only to the builders of slot $k$, during that slot. The auction assumes the fast finality planned for lean consensus, under which the set is fixed within the reference $T = 30$\,s of exposure (\cref{tab:choosing-t}). Under today's finality the delay would have to cover about 13 minutes, with the larger work and costs of the last row of \cref{tab:choosing-t} (\cref{sec:reorgs}).
  \item \textbf{Reveal.} The reveal phase begins once slot $k$'s block is final, so no key is published before the set of bids is fixed. Within the next $W$ slots, the bidder publishes its key. The contract checks it against $C$, checks the tag, decrypts the bid and returns $d$ and $B$. An invalid bid is excluded and forfeits $d$.
  \item \textbf{Forced opening.} A bid still sealed after the window is assigned to a solver, as in \cref{sec:bond}, drawn with the randomness of a later slot. A key matching $C$, from anyone, opens the bid, which is then checked as in a reveal and returns $B$; $d$ pays the solver. A failing tag, or a malformation claim that stands, excludes the bid and forfeits $d$ and $B$.
  \item \textbf{Settlement.} Once every bid is open or excluded, the contract sorts the valid bids by price, highest first, breaking ties by a hash of the sealed bid and the randomness of a later slot, fills them from $Q_k$ in that order, the last one partly, and charges every filled bid the lowest filled price. The rest of each escrow is refunded, and unsold units return to the seller. Later batches do not wait for this one.
\end{enumerate}

Settlement follows the reveal window at once if every bidder reveals. After a missed reveal, if the first solver assigned is honest and live and its moves are included in time, it waits from the assignment at most $D_s$ for an opening or a failing tag, $D_s + D_c$ for an unchallenged claim, and $D_s + D_c + 69\tau$ for a disputed one. With $D_s = D_c = 10$ minutes and $\tau = 2$ minutes, these are 10, 20 and 158 minutes. Each failed assignment can add as much again, and costs a staked solver its stake.

\Cref{tab:auction} compares this auction with commit--reveal under the same rules. In the auction, a valid, adequately escrowed bid with a well-formed puzzle remains binding after a missed reveal, including that of the last revealer, who sees the others' bids before its deadline. Withholding a malformed bid incurs an additional bond penalty, and an honest offline bidder still forfeits the reveal deposit. Commit--reveal can impose the same $d + B$ penalty on every missed reveal, but then an offline bidder pays it too. Against this, \leanseal costs bytes, the bidder's work, solver work and settlement delay after a missed reveal. Bids are private only for the delay: a bid that the builder leaves out of slot $k$ is not counted and opens anyway (\cref{sec:noninclusion}).

\begin{table}[t]
  \centering\small
  \caption{The auction under commit--reveal and under \leanseal, with the same escrow, admission slot and reveal window. Settlement times assume an honest, live first solver whose moves are included in time, and the deadlines of the text.}
  \label{tab:auction}
  \begin{tabularx}{\textwidth}{@{}>{\raggedright\arraybackslash}p{2.6cm}>{\raggedright\arraybackslash}X>{\raggedright\arraybackslash}X>{\raggedright\arraybackslash}X@{}}
    \toprule
    & \colhead{Commit--reveal, deposit $d$} & \colhead{Commit--reveal, deposit $d + B$} & \colhead{\leanseal, deposit $d$ and bond $B$} \\
    \midrule
    Bidder goes offline & bid lost; forfeits $d$ & bid lost; forfeits $d + B$ & a valid bid settles; forfeits $d$, which pays the solver \\
    Bidder withholds after seeing the market or others' bids & cancels for $d$ & cancels for $d + B$ & well-formed puzzle: a valid bid settles anyway; malformed: cancels for $d + B$ \\
    Added per bid & a 32-byte commitment & a 32-byte commitment & 152 to 512 bytes; about 21\,s of a laptop, ahead of time \\
    Solver work & none & none & 3.3 core-minutes per missed reveal \\
    Settlement & after $W$ slots & after $W$ slots & after the reveal window if all bids reveal; otherwise, from the first solver assignment, at most 10 minutes for an opening or failing tag, 20 for an unchallenged claim, and 158 for a disputed claim \\
    \bottomrule
  \end{tabularx}
\end{table}

\subsection{Other requirements}\label{sec:protocol}\label{sec:admission}\label{sec:reorgs}\label{sec:solvers}

\paragraph{Finality.}
The deadline $\torder$ is when the position can no longer change. Today an Ethereum block can be reorganized until it is finalized, about 13 minutes later, far beyond a 30-second lock; after a reorganization, a revealed transaction can be front-run in the replacement block. A protocol can accept this once a block has its attestations, so that a reorganization degrades only the affected transactions to today's exposure. With the faster finality planned for lean consensus, $\torder$ shrinks to seconds and a 30-second lock covers it with a wide margin.

\paragraph{Admission.}
A sealed transaction cannot be checked before it is opened, so, as in other encrypted mempools, its plaintext envelope is signed and checked as usual and its fee is charged whether or not the opened transaction is valid. The number of open puzzles per sender should be capped.

\paragraph{Solvers.}
Opening a withholding sender's transaction within about the delay needs a chip as fast as the assumed attacker; a laptop core takes 3.3 minutes for the reference puzzle. If every one of the roughly 265 transactions per slot in our mainnet sample had to be forced open, this would occupy about 660 attacker-speed circuits or 4,400 laptop cores; at a forced-opening rate of 1\%, about 44 cores, before redundancy and failed solves. With deposits in place, withholding should be rare.

A published key can be copied, and a withholding sender can publish its own key after the deadline, so a reward for the first matching key would go to copiers or to the sender; designated solvers paid per assignment avoid the first, and the second unless the sender controls the solver (\cref{sec:bond}). Whatever payment the sender or its coalition can recover reduces its penalty, so the cost of withholding is the forfeited deposit minus every such payment, and it is this difference that must deter withholding (\cref{sec:withholding}).

\paragraph{Size and gas.}
\label{sec:size}Sealing adds 152 bytes with stockpiled single-chain puzzles and 512 with the reference puzzle. In 20 consecutive mainnet blocks (22 October 2026; 5,300 transactions; reproducible with \texttt{bench/eth-sample.py}), the median transaction was 179 bytes, so sealing makes it 1.8 or 3.9 times larger. As calldata at 16 gas per nonzero byte the added bytes cost about 2,400 or 8,200 gas, against 21,000 for a plain transfer, and about 6,100 or 20,500 for a transaction that pays the calldata floor of EIP-7623, 40 gas per nonzero byte. The ordering needs only the 32-byte hash of each sealed transaction; the rest need only remain available until it is opened, for which a blob, retained for about 18 days, suffices.

\paragraph{Sharing puzzles.}\label{sec:sharing}
One puzzle can cover all of a sender's transactions that should open together. It cannot be shared with unrelated senders: everyone who encrypts under a key knows it. Sharing the cost of opening across senders is delay encryption, which hashes cannot provide when the setup is cheap (\cref{sec:limits}). A sender without sufficient hardware can instead outsource building to several services, and encrypt under a hash of all their keys, which stays secret if any one service is honest. Opening, however, needs every service's puzzle to be well formed, and the transaction carries all of them, so its size and the solvers' work grow with the number of services.
